\newpage
% =================================================
\section{Regras de Segurança do Firestore (Security Rules)}
% =================================================

\begin{explicacao}[Por que as Security Rules são obrigatórias]
Toda leitura via \texttt{onSnapshot} no frontend passa \textbf{diretamente} pelo Firestore SDK, \textbf{sem} passar pelas Cloud Functions. Portanto, as Cloud Functions são a camada de autorização para escrita, mas as \textit{Security Rules} são a camada de autorização para \textbf{leitura direta}. Sem rules corretas, qualquer usuário autenticado poderia ler qualquer coleção.
\end{explicacao}

\begin{regra}[Princípio: Deny-All por Padrão]
A regra raiz do arquivo \texttt{firestore.rules} deve ser:\\
\texttt{allow read, write: if false;}\\
Todas as exceções são declaradas explicitamente abaixo.
\end{regra}

\subsection{Regras por Coleção}

\small

\small
\setlength{\tabcolsep}{5pt}
\small
\setlength{\tabcolsep}{4pt}
\begin{longtable}{>{\raggedright\arraybackslash}p{0.32\textwidth} >{\raggedright\arraybackslash}p{0.62\textwidth}}
\toprule
\textbf{Coleção} & \textbf{Regra}\\
\midrule
\endfirsthead
\toprule
\textbf{Coleção} & \textbf{Regra} (continuação)\\
\midrule
\endhead
\bottomrule
\endfoot
\nolinkurl{Usuarios/\{uid\}} & Read: \texttt{request.auth.uid == uid} (usuário lê só o próprio documento). Write: negado ao cliente (só Cloud Function).\\
\nolinkurl{Usuarios/\{uid\}/Turmas/\{turmaId\}} & Read: \texttt{request.auth.uid == uid} (o aluno lê em tempo real apenas as turmas em que está matriculado). Write: negado ao cliente (gerenciada exclusivamente pelas Cloud Functions).\\
\texttt{Chefe\_Geral}, \texttt{Gestor\_Almoxarifado}, etc.~(coleções de papel) & Read: \texttt{request.auth.uid == docId} (cada usuário lê só o próprio documento de papel). Write: negado ao cliente.\\
\nolinkurl{Turma/\{turmaId\}} & Read: usuário é professor da turma (\texttt{resource.data.id\_professor == request.auth.uid}) OU está na sub-coleção \texttt{Alunos}. Write: negado ao cliente.\\
\nolinkurl{Turma/\{turmaId\}/Alunos} & Read: \texttt{request.auth.uid == resource.data.id\_aluno} OU membro ativo da mesma turma, apenas em projeção mínima sem matrícula/e-mail.\\
\nolinkurl{Turma/\{turmaId\}/Posts} & Read: membro da turma (professor ou aluno matriculado). Write: negado.\\
\nolinkurl{Turma/\{turmaId\}/Posts/\{postId\}/Comentarios} & Original moderado: somente autor, Professor responsável ou Chefe Geral com acesso ao contexto. Colegas recebem aviso por endpoint autorizado. Escrita direta negada.\\
\nolinkurl{Usuarios/\{uid\}/Notificacoes} & Read: \texttt{request.auth.uid == uid}. Write: negado ao cliente.\\
\nolinkurl{Resumo\_Reagente/\{resumoId\}} & Read: qualquer usuário com papel ativo no sistema (token de acesso válido). Write: negado ao cliente.\\
\nolinkurl{Resumo\_Reagente/\{resumoId\}/Especificacoes/\{specId\}} & Read: qualquer usuário com papel ativo no sistema (catálogo público interno). Write: negado ao cliente.\\
\texttt{Frasco\_Reagente} & Read: papel \texttt{roles} no token contém \texttt{Chefe\_Geral} OU \texttt{Gestor\_Almoxarifado} OU \texttt{Professor} OU \texttt{Bolsista} OU \texttt{Aluno}. Write: negado ao cliente.\\
\texttt{Emprestimo\_Reagente} & Read: \texttt{Chefe\_Geral} OU \texttt{Gestor\_Almoxarifado} OU o próprio usuário que retirou (\texttt{resource.data.id\_usuario\_retirou == request.auth.uid}). Write: negado.\\
\texttt{Bem\_Patrimonial} & Read: Chefe Geral, Gestor de Bens Patrimoniais ou Professor. Write: negado ao cliente.\\
\texttt{Requisicao\_Edicao\_Bem\_Patrimonial} & Read: o solicitante (\texttt{resource.data.id\_usuario\_solicitante == uid}) OU gestores de bens (papel no token). Write: negado.\\
\texttt{Requisicao\_Adicao\_Bem\_Patrimonial} & Idem acima.\\
\texttt{Locks\_Requisicao\_Patrimonio} & Read/Write: negado ao cliente (exclusivo de Cloud Function).\\
\texttt{Registro\_de\_Auditoria} & Read: \texttt{Chefe\_Geral} no token. Write: negado ao cliente.\\
\texttt{Roteiro\_Experimento} & Read: \texttt{resource.data.id\_professor\_upload == uid} OU \texttt{uid} pertence ao array \texttt{professores\_compartilhados}. Write: negado.\\
\texttt{Contador\_Codigo\_Frasco} & Read: \texttt{request.auth != null} (qualquer logado lê). Write: negado ao cliente (só Cloud Function atualiza o sequenciador).\\
\end{longtable}


\begin{regra}[Leitura de Custom Claims nas Security Rules]
O papel do usuário fica disponível nas Security Rules como \texttt{request.auth.token.roles}. Exemplo:\\
\texttt{function temPapel(papel) \{}\\
\texttt{\quad return papel in request.auth.token.roles;}\\
\texttt{\}}\\
Isso garante consistência com o mesmo Custom Claim usado pelas Cloud Functions (\texttt{resolverPapeisDoToken}).
\end{regra}


\subsection{Contrato de acesso para os fluxos detalhados}
\label{sec:acesso-detalhado}
Além do papel, validar o escopo atual. DP-D01 determina consulta de conta ativa em mutações de alto impacto conforme a seção 10; leituras e demais mutações podem confiar no JWT até renovação. Atualizar Custom Claims não invalida imediatamente o token emitido. Versão de permissões e reconciliação não constituem garantia universal já implementada. Preferência de papel no cliente só altera apresentação.

Leituras de membros, posts, comentários e PDFs de turma exigem professor responsável ou vínculo ativo; colegas recebem projeção mínima, sem e-mail/matrícula. Requisições são visíveis ao solicitante, gestor patrimonial e Chefe. Empréstimos completos são visíveis ao retirante, gestor vinculado ao almoxarifado e Chefe; demais leitores do catálogo recebem somente disponibilidade. Histórico operacional segue o mesmo escopo, e auditoria global fica com a chefia. Diretórios usados em seletores devem ter endpoints/projeções mínimas autorizadas, não leitura irrestrita de Usuarios.

Mutações de domínio e locks são exclusivas do servidor, inclusive marcar notificações lidas; não permitir delete de notificações para implementar Limpar tudo. Regras devem cobrir explicitamente Resumo\_Reagente/id/Especificacoes, Turma e suas subcoleções e Bem\_Patrimonial/id/Historico. Consultas collection-group exigem regras/índices próprios. Contador, unicidade e controles de papéis são técnicos; somente resposta mínima autorizada é exposta ao operador.

Storage valida dono, caminho, tipo e tamanho: foto inferior a 5 MiB, baixa inferior a 10 MiB e roteiro inferior a 15 MiB. Backend verifica arquivo existente e assinatura de conteúdo antes de referenciar; metadados não bastam para provar formato. Leitura de roteiro exige propriedade, compartilhamento ou post de turma acessível. Comprovante de baixa segue escopo patrimonial e não pode ser removido após vinculação. Não conceder leitura de todo Storage apenas porque o usuário está autenticado.

Estas são exigências da especificação; firestore.rules e storage.rules atuais ainda divergem, conforme AUDITORIA\_ATUALIZADA.md. A presença de deny-all genérico não anula permissões amplas de matches específicos.

\subsection{Moderação: autorização por documento e resposta filtrada}
Firestore Rules autorizam ou negam documentos completos; não fazem projeção nem substituição do campo texto. Não implementar ``texto vazio via Rules''. Para o modelo físico descrito, negar leitura direta de Comentarios e histórico a quem não pode conhecer o original; a listagem de comentários é fornecida por endpoint autenticado que valida vínculo com a turma e devolve somente aviso institucional para comentários moderados de colegas. Autor recebe original marcado; Professor responsável e Chefe Geral recebem conteúdo necessário à auditoria. Não retornar original em campos auxiliares, cache, notificações ou mensagens de erro. Consulta ampla cliente que possa incluir documentos proibidos deve ser rejeitada e substituída pelo endpoint, nunca filtrada apenas no frontend. Backend valida também autoria da edição, motivo e papel de moderação, status ativo da turma e registra histórico/auditoria.

Coleções Controle\_Papeis, Operacoes e Chaves\_Unicas são exclusivas de Admin SDK; leitura e escrita diretas de clientes são negadas. O deny-all atual já as protege, sem allow concorrente; matches explícitos constituem hardening documental. Esta seção descreve o alvo; a implementação atual de moderação permanece pendente.
